\documentclass[aspectratio=169,11pt]{beamer}

\usepackage{fontspec}
\usepackage{xeCJK}
\usepackage{graphicx}
\usepackage{booktabs}
\usepackage{amsmath}
\usepackage{url}
\usepackage{hyperref}

\usetheme{Madrid}
\usecolortheme{default}
\setbeamertemplate{navigation symbols}{}
\setbeamertemplate{footline}[frame number]

\setmainfont{Arial}
\setsansfont{Arial}
\IfFontExistsTF{Microsoft YaHei}{
  \setCJKmainfont{Microsoft YaHei}
  \setCJKsansfont{Microsoft YaHei}
}{
  \setCJKmainfont{SimSun}
  \setCJKsansfont{SimSun}
}

\graphicspath{{../figures/}}

\title{Predicting Loss Curve of LLM Pretraining}
\subtitle{Course Project Task 2}
\author{王渭臻，1900010767，本科生}
\institute{GitHub repository: \url{https://github.com/Tomori47/task2\_loss\_prediction}}
\date{\today}

\begin{document}

\begin{frame}
  \titlepage
\end{frame}

\begin{frame}{Problem Background}
  \begin{itemize}
    \item LLM 预训练成本很高，因此希望提前预测训练过程中的 loss curve。
    \item Scaling laws 研究模型规模、数据规模、训练计算量和 loss 的关系。
    \item Loss curve prediction 进一步关注整个训练过程，而不只是最终 loss。
    \item Learning rate schedule 会影响后期 loss 下降速度和最终效果。
  \end{itemize}
\end{frame}

\begin{frame}{Task Objective}
  \begin{itemize}
    \item 课程 Task 2 要求：使用 cosine LRS 的 loss curve 拟合模型。
    \item 主测试：预测 WSD LRS 下的 loss curve。
    \item 附加测试：在 811 LRS 上评估泛化效果。
    \item 评价指标：MAE、MSE、RMSE、MAPE。
  \end{itemize}
\end{frame}

\begin{frame}{Data Description}
  \begin{itemize}
    \item 数据文件：\texttt{data/gpt\_loss+lrs.pkl}
    \item 三条曲线：\texttt{cosine\_rope}、\texttt{wsd\_rope}、\texttt{811\_rope}
    \item 每条曲线包含：\texttt{step}、\texttt{Metrics/loss}、\texttt{lr}
  \end{itemize}
  \vspace{0.2cm}
  \begin{columns}
    \column{0.49\textwidth}
    \includegraphics[width=\textwidth]{\detokenize{data_preview_loss.png}}
    \column{0.49\textwidth}
    \includegraphics[width=\textwidth]{\detokenize{data_preview_lr.png}}
  \end{columns}
\end{frame}

\begin{frame}{Experimental Setting}
  \begin{table}
    \centering
    \begin{tabular}{ll}
      \toprule
      Split & Curve \\
      \midrule
      Fit & \texttt{M:100M\_gpt\_D:20B\_scheduler:cosine\_rope} \\
      Main test & \texttt{M:100M\_gpt\_D:20B\_scheduler:wsd\_rope} \\
      Additional test & \texttt{M:100M\_gpt\_D:20B\_scheduler:811\_rope} \\
      \bottomrule
    \end{tabular}
  \end{table}
  \vspace{0.3cm}
  \begin{itemize}
    \item 所有方法使用相同的 fit/test 划分。
    \item 目标是检查从 cosine schedule 到 WSD schedule 的预测能力。
  \end{itemize}
\end{frame}

\begin{frame}{Baseline 1: Tissue et al.}
  \begin{block}{Model}
    \[
      L(s) = L_0 + A\cdot S_1(s)^{-\alpha} - C\cdot S_2(s)
    \]
  \end{block}
  \begin{itemize}
    \item $S_1(s)$：累计学习率。
    \item $S_2(s)$：learning rate annealing 修正项。
    \item 本项目中固定 \texttt{lambda\_decay = 0.99}。
    \item 直觉：loss 随累计训练强度下降，同时由学习率衰减项修正。
  \end{itemize}
\end{frame}

\begin{frame}{Baseline 2: Simplified Multi-Power Law}
  \begin{block}{Simplified MPL}
    \[
      L(t) = L_0 + A\cdot(S_1(t)+\epsilon)^{-\alpha} - B\cdot D(t)
    \]
  \end{block}
  \[
    D(t) = \sum_{k=1}^{t} \max(lr_{k-1} - lr_k, 0)
           \cdot(S_{\text{tail}}(k,t)+\epsilon)
  \]
  \begin{itemize}
    \item 这是简化实现，没有完全复现 Luo et al. 的所有训练细节。
    \item 保留两个核心思想：累计学习率项和学习率衰减修正项。
  \end{itemize}
\end{frame}

\begin{frame}{Our Method: Weighted Tissue}
  \begin{itemize}
    \item 预测公式仍然使用 Tissue 形式。
    \item 改动只在拟合目标：对后期 step 给予更高权重。
  \end{itemize}
  \begin{block}{Weighted Least Squares}
    \[
      \min \sum_s w_s \left(L_{\text{true}}(s)-L_{\text{pred}}(s)\right)^2
    \]
    \[
      w_s = 1 + \gamma\cdot\frac{s}{T}, \quad \gamma = 2.0
    \]
  \end{block}
  \begin{itemize}
    \item 直觉：后期 loss 和最终 loss 更重要，因此拟合时更重视后期曲线。
  \end{itemize}
\end{frame}

\begin{frame}{Results on Cosine Fit}
  \begin{table}
    \centering
    \begin{tabular}{lccc}
      \toprule
      Method & MAE & RMSE & MAPE \\
      \midrule
      Tissue & 0.07730 & 0.15378 & 2.40\% \\
      Simplified MPL & 0.07730 & 0.15378 & 2.40\% \\
      Weighted Tissue & 0.07200 & 0.15463 & 2.20\% \\
      \bottomrule
    \end{tabular}
  \end{table}
  \begin{itemize}
    \item 三种方法都能拟合 cosine loss 的主要趋势。
    \item Weighted Tissue 的 MAE 和 MAPE 更低，但 RMSE 略高。
  \end{itemize}
\end{frame}

\begin{frame}{Results on WSD Prediction}
  \begin{itemize}
    \item 主测试：在 WSD LRS 上预测 loss curve。
    \item 下图比较 ground truth、Tissue、simplified MPL 和 Weighted Tissue。
  \end{itemize}
  \centering
  \includegraphics[width=0.86\textwidth]{\detokenize{method_comparison_wsd.png}}
\end{frame}

\begin{frame}{Weighted Tissue on WSD}
  \begin{columns}
    \column{0.45\textwidth}
    \begin{itemize}
      \item MAE = 0.13999
      \item RMSE = 0.19436
      \item MAPE = 4.64\%
      \item 相比 baseline 有小幅改善。
    \end{itemize}
    \column{0.55\textwidth}
    \includegraphics[width=\textwidth]{\detokenize{our_method_predict_wsd.png}}
  \end{columns}
\end{frame}

\begin{frame}{Quantitative Comparison}
  \begin{table}
    \centering
    \begin{tabular}{lccc}
      \toprule
      Method & MAE & RMSE & MAPE \\
      \midrule
      Tissue & 0.15364 & 0.20276 & 5.14\% \\
      Simplified MPL & 0.15364 & 0.20275 & 5.14\% \\
      Weighted Tissue & 0.13999 & 0.19436 & 4.64\% \\
      \bottomrule
    \end{tabular}
  \end{table}
  \begin{itemize}
    \item Weighted Tissue 将 WSD MAE 从约 0.15364 降到 0.13999。
    \item Weighted Tissue 将 WSD MAPE 从约 5.14\% 降到 4.64\%。
  \end{itemize}
\end{frame}

\begin{frame}{Analysis}
  \begin{itemize}
    \item Tissue 和 simplified MPL 的衰减修正项在当前数据中几乎收敛到 0。
    \item 说明只用一条 cosine 曲线拟合时，模型主要依赖累计学习率 power-law 项。
    \item Weighted Tissue 没有改变模型形式，只改变拟合权重。
    \item 通过强调后期 loss，Weighted Tissue 在 WSD 预测上获得小幅改善。
  \end{itemize}
\end{frame}

\begin{frame}{Limitations}
  \begin{itemize}
    \item 只使用了三条 loss curves，数据规模较小。
    \item Simplified MPL 没有完全复现 Luo et al. 的完整模型。
    \item 没有进行大规模超参数搜索。
    \item Weighted Tissue 只使用固定的 $\gamma=2.0$。
    \item 结果主要作为低成本课程复现实验。
  \end{itemize}
\end{frame}

\begin{frame}{Conclusion}
  \begin{itemize}
    \item 完成 Tissue baseline。
    \item 完成 simplified Multi-Power Law baseline。
    \item 提出并实现 Weighted Tissue fitting。
    \item 在 WSD 主测试上，Weighted Tissue 相比 baseline 有小幅改善。
    \item 已生成代码、CSV 指标、图像、README 和 slides。
  \end{itemize}
\end{frame}

\begin{frame}{Code and Contributions}
  \begin{itemize}
    \item GitHub repository: \url{https://github.com/Tomori47/task2\_loss\_prediction}
    \item Team members: 王渭臻，1900010767，本科生
    \item 关键输出：
      \begin{itemize}
        \item \texttt{results/metrics\_summary.csv}
        \item \texttt{figures/method\_comparison\_wsd.png}
        \item \texttt{slides/task2\_loss\_prediction\_slides.tex}
      \end{itemize}
  \end{itemize}
\end{frame}

\begin{frame}{References}
  \begin{itemize}
    \item Tissue et al., 2024. \textit{Scaling Law with Learning Rate Annealing}.
    \item Luo et al., 2024/2025. \textit{Multi-Power Law Scaling for Loss Curve Prediction}.
    \item Course Task 2 handout: \textit{Predicting Loss Curve of LLM Pretraining}.
  \end{itemize}
\end{frame}

\end{document}
